\scanpage{14}
\begingroup\makeatletter\tagsleft@true\makeatother
\[
\tau_K\circ u_\xi^K=\tau_n\tag*{(1)}
\]
et la propriété de transitivité pour les applications admissibles
\[
\begin{aligned}[t]
&\forall\,\xi\in K,\ \eta\in J_n\quad(n=\tau_K(\xi))\ \text{on a, si }n'=\tau_n(\eta)\\
&\qquad\deleted{\unsure{\tau|u}}\,u_\xi^K\circ u_\eta^n=u_\zeta^K,\quad\text{où }\zeta=u_\xi^K(\eta),
\end{aligned}\tag*{(2)}
\]
\endgroup
enfin

\deleted{(3) Si $\xi\in K,\ \eta\in K$, alors}

La catégorie dont les objets sont les objets admissibles de $\widehat M$, les flèches les \unsure{Isom|isom|mono|hom} entre eux, est équivalente : celle des ens ordonnés $M$-épinglés.

D'autre part, les hom quelconques \inserted{dans $\widehat C$ entre tels objets} se réconstituent également, ce sont les applictions \uline{strictement} \uline{croissantes}
\[
K\xrightarrow{f}L
\]
telles que $\left\{\begin{aligned}
&\text{a) }\tau_L\circ f=\tau_K\\
&\text{b) pour tt }\xi\in K,\ \inserted{\text{posant }\eta=f(\xi)\in L,}\ \text{on ait}
\end{aligned}\right.$ commutativité \unsure{du|de}
\diagram{\begin{tikzcd}[ampersand replacement=\&,column sep=small,row sep=large]
K_{\leq\xi}\arrow[rr,"{f|K_{\leq\xi}}"] \& \& {\deleted{K}L}_{\leq\eta}\\
\& J_n\arrow[ul,"\sim"']\arrow[ur,"\sim"] \& \& \text{où }n=\tau_K(\xi)=\tau_{\deleted{\unsure{K|\eta}}L}(\eta),
\end{tikzcd}}
% archive sheet number: 13
